\section{Conclusiones y uso para crédito y seguro}\label{sec:conclusiones}

Construimos un modelo del rendimiento municipal de la cebada de temporal con datos públicos y lo validamos como si
cada municipio fuera un lugar nuevo. La temperatura máxima de julio y agosto, corregida con el kriging de los
residuos, predice el nivel municipal con un RMSE ponderado de 0.293\,t/ha con datos cerca y de 0.438\,t/ha sin datos a
menos de 50\,km. El kriging de vecinos tiene 0.309 y 0.566. El suelo no agrega información detectable y el clima
mensual del año no pronostica la cosecha.

Para FIRA proponemos tres usos.
\begin{enumerate}
\item \textbf{Priorizar zonas para crédito.} El rendimiento esperado y su intervalo p10 a p90 ordenan los municipios
por su potencial. Los de confianza alta tienen una mediana de 2.53\,t/ha. Los de confianza media sirven para ampliar
la cobertura a municipios vecinos sin historia propia.
\item \textbf{Tarifar el seguro.} La probabilidad de quedar bajo 2\,t/ha en un año cualquiera resume el riesgo de cada
municipio. Entre los municipios de confianza alta va de 0.00 en Lázaro Cárdenas, Tlaxcala, a 0.99 en San
Agustín Tlaxiaca, Hidalgo. Como sobreestima un poco la
frecuencia observada, es una cota prudente. Para el seguro conviene usar el rendimiento sobre sembrada, que incluye la
pérdida por siniestro.
\item \textbf{Exigir verificación donde la confianza es baja.} En los 77 municipios de confianza baja el valor solo
refleja el clima. Antes de financiar ahí hace falta confirmar que se siembra cebada y que el suelo es agrícola.
\end{enumerate}

El modelo predice el promedio de un municipio. El rendimiento de un productor depende además de su manejo, y las
etiquetas del régimen A del reto quedan 1.71\,t/ha arriba del modelo. Para estimar el rendimiento de una parcela
harán falta datos de manejo y rendimientos medidos con un mismo criterio en todos los municipios.
